\chapterimage{chapter_pointers.jpg} % Chapter heading image

\chapter{Visão Geral da linguagem C}

Se você chegou até aqui, significa que já passou pelo batismo inicial de conceitos fundamentais. Excelente! Agora é hora de colocar as mãos na massa e entender como a linguagem C realmente funciona.

\section{Origens da linguagem}

C é uma linguagem antiga. Muito antiga. Estamos falando de 1972, quando o planeta ainda não era conectado pela internet como hoje. Naquela época, Dennis Ritchie, pesquisador dos Bell Labs, precisava de uma linguagem que fosse poderosa o suficiente para manipular hardware, mas também legível para humanos.

\begin{quote}
\textit{Por que uma linguagem tão antiga ainda é relevante nos dias de hoje?}
\end{quote}

Ótima pergunta! A resposta é simples: C foi revolucionária. Ela oferecia o melhor dos dois mundos: a proximidade com o hardware das linguagens de baixo nível e a legibilidade das linguagens de alto nível. Sistema operacional Unix, que mudou o mundo da computação, foi escrito em C. Até hoje, o kernel do Linux é predominantly escrito em C.

Portanto, se você quer entender como os computadores funcionam por baixo, C é a porta de entrada. É como aprender latim para entender as raízes das línguas modernas, mas, felizmente, bem mais útil e com uma comunidade ativa.

\section{Porque aprender C?}

Você pode estar se perguntando: existem linguagens mais modernas e fáceis por aí. Por que diabos aprender C?

\begin{enumerate}
\item \textbf{Entendimento de Fundamentos}: C te força a pensar sobre memória, ponteiros e estrutura de dados. Essas são habilidades que você levará para qualquer linguagem.
\item \textbf{Performance}: Código bem escrito em C é extraordinariamente rápido. Sistemas embarcados, bancos de dados e servidores críticos usam C porque ela entrega performance.
\item \textbf{Portabilidade}: Um programa em C compilado funciona em praticamente qualquer plataforma. Windows, Linux, macOS, sistemas embarcados, tudo funciona com C.
\item \textbf{Oportunidades Profissionais}: Muitas empresas ainda usam C em producao. Conhecer C te abre portas em áreas como sistemas operacionais, compiladores, databases e infraestrutura.
\item \textbf{Simplicidade}: Comparada a linguagens modernas como C++, C é simples. Há poucas abstrações entre você e o que realmente está acontecendo na máquina.
\end{enumerate}

\section{Onde C é usado?}

C está em todos os lugares, mesmo que você não veja diretamente.

\subsection{Sistemas Operacionais}

Linux, Windows (em partes), macOS, Android e iOS têm componentes críticos escritos em C. O kernel do Linux, responsável pelo gerenciamento de recursos do sistema, é principalmente C.

\subsection{Bancos de Dados}

SQLite, PostgreSQL e MySQL têm suas implementações core em C. Se você usa um banco de dados, há grandes chances de estar usando código C por baixo.

\subsection{Sistemas Embarcados}

Microcontroladores, IoT devices, smartwatches - tudo isso frequentemente roda C. Quando recursos de memória são limitados, C é a escolha obvia.

\subsection{Compiladores e Ferramentas}

GCC (GNU Compiler Collection), LLVM e muitas outras ferramentas de desenvolvimento são escritas em C ou C++. Python, Node.js e outras linguagens têm parsers e engines em C.

\subsection{Infraestrutura e Performance-Critical Software}

Redis, Nginx, Apache e outros softwares de alta performance usam C.

\section{Estrutura de um programa em C}

Todo programa em C segue uma estrutura básica. Vejamos um exemplo simples:

\begin{ccode}
  #include <stdio.h>

  int main(void)
  {
    printf("Olá, Mundo!\n");
    return 0;
  }
\end{ccode}

Parece simples, não é? Mas cada linha tem um propósito. Vamos dissecar:

\begin{itemize}
  \item \textbf{\#include <stdio.h>}: Esta é uma \textit{diretiva de pré-processador}. Ela diz ao compilador para incluir a biblioteca padrão de entrada/saída (Standard Input/Output). A função \textit{printf} vem dessa biblioteca.

  \item \textbf{int main(void)}: Esta é a \textit{função principal}. Todo programa em C deve ter uma função \textit{main}. O compilador começa a execução daqui. O \textit{int} indica que a função retorna um inteiro. O \textit{void} significa que ela não recebe parâmetros.

  \item \textbf{printf(...)}: Esta função imprime texto na tela. A sequência \textit{\textbackslash n} é um caractere de quebra de linha.

  \item \textbf{return 0;}: Retorna o valor zero para o sistema operacional, indicando que o programa terminou com sucesso.
\end{itemize}

\section{Meu primeiro programa em C}

Vamos criar e executar seu primeiro programa. Para isso, você precisa de:

\begin{enumerate}
  \item Um \textit{editor de texto} (qualquer um serve: gedit, vim, VS Code, etc)
  \item Um \textit{compilador C} (GCC é recomendado e gratuito)
\end{enumerate}

Crie um arquivo chamado \textbf{hello.c} com o seguinte conteúdo:

\begin{ccode}
  #include <stdio.h>

  int main(void)
  {
    printf("Olá, eu sou um programador C!\n");
    return 0;
  }
\end{ccode}

Salve o arquivo. Parabéns, você acabou de criar seu primeiro programa em C!

\section{Compilando meu programa}

Agora vamos transformar esse código em um programa executável. Abra seu terminal e navegue até a pasta onde salvou o arquivo \textbf{hello.c}.

\begin{quote}
\textit{Terminal? Isso é aquela tela preta assustadora?}
\end{quote}

Sim, mas não se preocupe. A linha de comando é sua amiga, especialmente para programadores.

Digite o seguinte comando:

\begin{ccode}
  gcc hello.c -o hello
\end{ccode}

O que aconteceu? Vamos entender:

\begin{itemize}
  \item \textbf{gcc}: O compilador GNU C.
  \item \textbf{hello.c}: O arquivo de código fonte que queremos compilar.
  \item \textbf{-o hello}: A flag \textit{-o} diz ao compilador qual nome dar ao executável de saída. Neste caso, \textit{hello}.
\end{itemize}

Se tudo correr bem, você verá apenas o prompt retornando. Isso é bom! Significa que não houve erros. Se houvesse, o compilador teria mostrado mensagens de erro.

\section{Executando}

Agora vamos executar o programa que compilamos. No terminal, digite:

\begin{ccode}
  ./hello
\end{ccode}

O ponto e barra \textit{./} significa "execute o arquivo da pasta atual". Se você estiver no Windows, pode simplesmente digitar \textit{hello} sem o ponto e barra.

Resultado esperado:

\begin{quote}
\textit{Olá, eu sou um programador C!}
\end{quote}

Parabéns! Você acabou de escrever, compilar e executar seu primeiro programa em C. Tira uma foto disso aí para mostrar pros amigos.

\section{Entendendo meu programa}

Vamos revisitar o código e entender o que realmente aconteceu:

\begin{ccode}
  #include <stdio.h>    // 1. Inclui biblioteca

  int main(void)        // 2. Função principal
  {                     // 3. Abre bloco
    printf("Olá, eu sou um programador C!\n");  // 4. Imprime texto
    return 0;           // 5. Retorna sucesso
  }                     // 6. Fecha bloco
\end{ccode}

\begin{enumerate}
  \item O \textit{\#include <stdio.h>} carrega a biblioteca padrão que contém a função \textit{printf}.
  \item \textit{main} é o ponto de entrada do programa. Quando você executa o programa, o sistema operacional chamará esta função.
  \item As chaves \textit{\{\}} delimitam o corpo da função.
  \item \textit{printf} escreve texto na saída padrão (geralmente a tela).
  \item \textit{return 0} indica que o programa terminou com sucesso (código de saída 0).
  \item A chave de fechamento marca o fim da função.
\end{enumerate}

O fluxo de execução é linear: de cima para baixo. O programa executa linha por linha até encontrar um \textit{return}.

Agora você tem um entendimento básico de como um programa em C funciona. No próximo capítulo, vamos explorar variáveis e tipos de dados.

